% !TeX program = xelatex
% ============================================================
% 2026 高教社杯全国大学生数学建模竞赛 B 题
% 问题一：有界测向误差下的交会定位区域
%
% 本文件按 CUMCM-LaTeX-Template/template.tex 的写法组织，
% 保留模板标题、交叉引用、定理环境、图表和 Step 求解结构。
% 本文件不含摘要、目录、参考文献、附录和后续问题。
%
% 两个图位对应的文件为：
%   figures/F-Q1-01.pdf
%   figures/F-Q1-02.pdf
% ============================================================

\PassOptionsToPackage{quiet}{xeCJK}
\documentclass[withoutpreface,notoc]{cumcmthesis}

% 与模板一致：正文中的加粗文字使用黑体。
\RenewDocumentCommand{\textbf}{m}{{\heiti #1}}

% 补充模板未显式列出的排版能力；模板已经加载的宏包不会重复加载。
\usepackage{amsmath,amssymb,bm}
\usepackage{array}
\usepackage{tabularx}
\usepackage{booktabs}
\usepackage{etoolbox}
\usepackage{enumitem}
\usepackage{graphicx}

% 表格使用模板约定的小五号字。
\BeforeBeginEnvironment{tabular}{\zihao{-5}}
\BeforeBeginEnvironment{tabularx}{\zihao{-5}}

\newcolumntype{L}{>{\raggedright\arraybackslash}X}
\newcolumntype{C}{>{\centering\arraybackslash}X}

% 常用记号。
\newcommand{\bs}{\boldsymbol{S}}
\newcommand{\bp}{\boldsymbol{p}}
\newcommand{\bq}{\boldsymbol{q}}
\newcommand{\bd}{\boldsymbol{d}}
\newcommand{\bu}{\boldsymbol{u}}
\newcommand{\bv}{\boldsymbol{v}}
\newcommand{\ba}{\boldsymbol{a}}
\newcommand{\bb}{\boldsymbol{b}}
\newcommand{\bA}{\boldsymbol{A}}
\newcommand{\bB}{\boldsymbol{B}}
\newcommand{\bC}{\boldsymbol{C}}
\newcommand{\bO}{\boldsymbol{O}}
\newcommand{\bzero}{\boldsymbol{0}}
\newcommand{\Ball}{\mathcal{B}}
\newcommand{\T}{\mathrm{T}}
\newcommand{\R}{\mathbb{R}}
\DeclareMathOperator{\diam}{diam}
\DeclareMathOperator{\conv}{conv}
\DeclareMathOperator{\dist}{dist}

% 有图时直接插入 PDF，无图时显示固定高度的图位。
% 这样可以在正式配图尚未完成时检查版式和页数。
\newcommand{\FigureSlot}[4]{%
  \begin{figure}[htbp]
    \centering
    \IfFileExists{figures/#2}{%
      \includegraphics[width=#1,height=0.140\textheight,keepaspectratio]{figures/#2}%
    }{%
      \fbox{%
        \begin{minipage}[c][0.080\textheight][c]{0.80\textwidth}
          \centering
          \textbf{图位：#3}\\[0.25em]
          将 PDF 命名为 \texttt{#2}，放入 \texttt{figures/} 目录。
        \end{minipage}%
      }%
    }%
    \caption{#3}
    \label{#4}
  \end{figure}
}

\setlist[itemize]{leftmargin=2em,itemsep=0.08em,topsep=0.18em}
\setlist[enumerate]{leftmargin=2em,itemsep=0.08em,topsep=0.18em}

\begin{document}

% 论文总标题由总文档统一生成；第一问文件只提供问题一的章节正文。

\section{问题一的模型建立与求解}

\subsection{问题分析}

每个检测点的示向度误差绝对值不超过 $1^\circ$。因此，干扰源不能简单等同于两条中心
示向线的交点，而应位于以检测点为顶点、以实测示向度为角平分线、半角为
$\alpha=1^\circ$ 的前向楔形内。对同一干扰源的全部前向楔形求交，得到题目要求的
纯测向交会区域 $P$。由于每个楔形都是两个闭半平面的交，$P$ 是闭凸集。求解时必须
先判断 $P$ 是否为空、是否有界；只有在 $P$ 为非空有界凸多边形时，区域直径才有
有限值。随后还需判断是否存在直径等于 $D$ 的覆盖圆。目标圆域、接收距离不超过
$1500$ m 和 $5$ m 过近条件作为物理审计条件，不替代纯测向交会区域。

\subsection{模型假设}

本文假设：同一干扰源的示向度观测已正确关联；检测点坐标准确，各点与干扰源近似共面，
示向度误差绝对值不超过 $1^\circ$；暂不考虑地形遮挡、多径传播和接收机内部延迟；
目标圆域以坐标原点 $\bO$ 为圆心，该假设只影响物理约束审计。

\subsection{前向楔形与半平面模型}

记第 $i$ 个检测点为 $\bs_i=(x_i,y_i)$，其示向度为 $\theta_i$。在
$\alpha=1^\circ$ 的误差半角下，第 $i$ 个前向楔形的下、上边界单位方向向量为
\begin{equation}
  \bu_i^-=
  \begin{bmatrix}
    \cos(\theta_i-\alpha)\\
    \sin(\theta_i-\alpha)
  \end{bmatrix},
  \qquad
  \bu_i^+=
  \begin{bmatrix}
    \cos(\theta_i+\alpha)\\
    \sin(\theta_i+\alpha)
  \end{bmatrix}.
  \label{eq:boundary-vectors}
\end{equation}

对二维向量 $\ba,\bb$，定义叉积 $[\ba,\bb]=a_xb_y-a_yb_x$。候选位置
$\bp=(x,y)$ 属于第 $i$ 个前向楔形，当且仅当
\begin{equation}
  [\bu_i^-,\bp-\bs_i]\ge 0,
  \qquad
  [\bu_i^+,\bp-\bs_i]\le 0.
  \label{eq:wedge}
\end{equation}

式\cref{eq:wedge}直接用单位向量判断点与楔形边界的相对位置，不需要计算
$\tan\theta_i$，因此不会在 $90^\circ$、$270^\circ$ 等特殊角度或角度跨越
$0^\circ$ 时出现数值奇异。将叉积展开后，每条示向度观测给出两个半平面约束
$\ba_k^\T\bp\le b_k$。全部 $2n$ 个约束的交集为
\begin{equation}
  P=\bigcap_{k=1}^{2n}\left\{\bp\in\R^2:\ba_k^\T\bp\le b_k\right\}.
  \label{eq:intersection}
\end{equation}

式\cref{eq:intersection}是问题一的核心模型。若约束互相矛盾，则
$P=\varnothing$；若方向约束不能限制区域远端，则 $P$ 无界；这两种情况下均不能
输出有限的区域直径。其几何含义和直径判定如\cref{fig:q1-main}所示。

\FigureSlot{0.82\textwidth}{F-Q1-01.pdf}
  {两站前向楔形交会区域、直径与同直径圆判定}{fig:q1-main}

\subsection{区域状态、顶点与直径}

用二维增量法判断\cref{eq:intersection}的可行性：若候选点违反新约束，则把候选点
限制在新边界直线上，并将已有约束化为关于该直线参数的一维区间；区间交为空时
$P$ 不可行。确认 $P$ 非空后，用衰退锥 $C_P=\{\bd:\bA\bd\le\bzero\}$ 判断
有界性：
\begin{equation}
  C_P=\{\bd:\bA\bd\le\bzero\},
  \qquad
   P\ \text{有界}\Longleftrightarrow C_P=\{\bzero\}.
  \label{eq:bounded}
\end{equation}

只需检查 $d_x=\pm1$ 和 $d_y=\pm1$ 四类归一化方向；任意非零衰退方向按正比例
缩放后必属于其中一类，因此该判定是完备的。区域有界时，枚举非平行边界线的交点，
删除不满足约束的交点，再用 Andrew 单调链算法得到逆时针凸包顶点
$V(P)=\{\bv_0,\bv_1,\ldots,\bv_{h-1}\}$。由凸集的极值性质，区域直径满足
\begin{equation}
  D=\diam(P)
  =\max_{\bp,\bq\in P}\|\bp-\bq\|_2
  =\max_{0\le i,j<h}\|\bv_i-\bv_j\|_2.
  \label{eq:diameter}
\end{equation}

由紧凸多边形的极值性质，可用旋转卡壳算法在 $O(h)$ 时间内求出一组最远点
$\bA,\bB$，再用全部顶点对的 $O(h^2)$ 枚举复核。两种算法结果一致时才进入
覆盖判定。

\subsection{同直径圆覆盖判定}

设最远点对满足 $D=\|\bA-\bB\|_2$。若半径为 $D/2$、圆心为 $o$ 的圆盘覆盖
$P$，则三角不等式给出
$D=\|\bA-\bB\|_2\le\|\bA-o\|_2+\|o-\bB\|_2\le D$。两端相等要求每一步都取
等号，因此 $\bA,o,\bB$ 共线，且唯一可能的圆心为最远点对的中点
$\bC=(\bA+\bB)/2$。

所以只需检查所有凸包顶点是否落在以 $\bA\bB$ 为直径的圆盘中：
\begin{equation}
  (\bv-\bA)^\T(\bv-\bB)\le 0,
  \qquad \forall \bv\in V(P).
  \label{eq:cover-dot}
\end{equation}

式\cref{eq:cover-dot}成立时，存在直径为 $D$ 的覆盖圆；否则不存在。区域直径和
最小覆盖圆直径不是同一个概念，\cref{fig:q1-counterexample}给出了不能省略该判定
的反例。

\FigureSlot{0.76\textwidth}{F-Q1-02.pdf}
  {等边三角形区域直径与最小覆盖圆直径的差异}{fig:q1-counterexample}

\subsection{模型求解与物理约束审计}

主模型的求解按以下步骤进行：
\begin{description}
  \item[Step 1：] 读取坐标和示向度，按式\cref{eq:boundary-vectors}构造楔形边界。
  \item[Step 2：] 按式\cref{eq:wedge}和式\cref{eq:intersection}形成半平面约束，
        判断区域的可行性和有界性。
  \item[Step 3：] 筛选边界交点，提取逆时针凸包 $V(P)$。
  \item[Step 4：] 求区域直径并进行顶点对复核，按式\cref{eq:cover-dot}判断
        同直径圆是否存在。
  \item[Step 5：] 审计物理约束并完成误差半角灵敏度分析。
\end{description}

物理约束审计包括
\begin{equation}
  \max_{\bv\in V(P)}\|\bv\|_2\le 1800,
  \qquad
  \max_{\bv\in V(P)}\|\bv-\bs_i\|_2\le 1500,
  \qquad
  \min_{\bp\in P}\|\bp-\bs_i\|_2>5.
  \label{eq:audit}
\end{equation}

前两项可在凸多边形顶点上检查，第三项需要计算点到多边形的最短距离。若三项均不活跃，
则保留 $P$；否则在 $P$ 上叠加目标圆域和接收圆盘裁剪，并排除过近圆盘内部。该修正
区域只作为补充审计结果，不替代式\cref{eq:intersection}的纯测向交会多边形。

\subsection{算例结果与检验}

题面没有提供现场坐标和示向度，因此使用合成算例检验模型，以下结果不表示实际干扰源
位置。取
$\bs_1=(-100,-100)$、$\bs_2=(100,-100)$，示向度分别为 $45^\circ$ 和
$135^\circ$。旋转卡壳与顶点对枚举给出一致结果
\begin{equation}
  D=100(\tan46^\circ-\tan44^\circ)
  =6.984153898\ \mathrm{m}.
  \label{eq:two-station-diameter}
\end{equation}

最远点对的中点圆半径为 $D/2=3.492076949$ m。逐点检查后，
$P\subseteq\Ball(\bC,D/2)$ 成立，因此该算例存在同直径覆盖圆。

再取边长为 $36$ m 的等边三角形作为反例。其区域直径为 $36$ m，但最小覆盖圆
直径为 $24\sqrt{3}\approx41.56921938$ m。以任意一条边为直径的圆都不能覆盖
第三个顶点，说明同直径圆判定是独立必要步骤。

数值检验中，$3\times4$ 矩形直径的计算误差为 $0$，旋转平移后的直径变化为
$2.49\times10^{-14}$ m，固定随机种子的 $120$ 组凸包上两种算法结果一致。

表\cref{tab:sensitivity}给出误差半角对双站算例直径的影响。随着 $\alpha$ 由
$0.50^\circ$ 增至 $1.50^\circ$，楔形按集合包含关系扩张，区域直径严格增大，
与测向误差界的几何含义一致。

\begin{table}[H]
  \centering
  \renewcommand{\arraystretch}{1.05}
  \caption{误差半角对双站算例区域直径的影响}
  \label{tab:sensitivity}
  \begin{tabular}{lccccc}
    \toprule
    误差半角 $/^\circ$ & 0.50 & 0.75 & 1.00 & 1.25 & 1.50\\
    \midrule
    区域直径 / m & 3.491 & 5.237 & 6.984 & 8.732 & 10.482\\
    \bottomrule
  \end{tabular}
\end{table}

\end{document}
